\chapter{System overview}
\label{ch:overview}

\section{Project goal}
FPGA-Neural implements a \textbf{reusable Neural Network Engine in FPGA hardware}.
The whole is made of three elements: the FPGA, which is the actual accelerator; a
dedicated RAM physically associated with the FPGA and not shared with the host; and a
host interface independent of the operating system, initially SPI (with possible
future extension to Dual~SPI).

The founding principle is the separation between who \emph{executes} the computation
and who \emph{uses} it: the neural network computation happens entirely inside the
FPGA, while the host system only provides configuration, network parameters, input
data, control and result readback. The host is not part of the computational datapath.
Possible host systems include Linux SoCs, Raspberry~Pi-like systems, ESP32,
microcontrollers and development PCs: the same engine architecture must be usable in
completely different systems.

\begin{center}
\begin{tikzpicture}[font=\footnotesize,node distance=8mm]
  \node[fnblockD,minimum width=42mm,minimum height=20mm] (host){\textbf{HOST}\\[2pt]
        {\scriptsize Configuration}\\{\scriptsize Training}\\{\scriptsize Control}};
  \node[fnblockT,below=14mm of host,minimum width=42mm,minimum height=20mm] (fpga)
        {\textbf{FPGA}\\[2pt]{\scriptsize Neural Network Engine}\\{\scriptsize Compute / Control}};
  \node[fnblock,below=14mm of fpga,minimum width=42mm,minimum height=13mm] (ram)
        {\textbf{Dedicated RAM}\\{\scriptsize weights / bias / buffers}};
  \draw[fnbus] (host) -- node[fnlbl,right]{SPI / Dual SPI} (fpga);
  \draw[fnbus] (fpga) -- node[fnlbl,right]{parallel bus} (ram);
\end{tikzpicture}
\end{center}

\section{Hardware configuration versus network configuration}
The project draws a precise distinction between the accelerator's \textbf{hardware
architecture} and the \textbf{neural network parameters}.

The physical architecture of the engine is defined at FPGA synthesis and
implementation time. Typical hardware parameters are \code{N\_INPUTS},
\code{N\_NEURONS}, \code{N\_LAYERS}, \code{PARALLEL}, \code{DATA\_WIDTH},
\code{ACC\_WIDTH}: they are Verilog parameters resolved at synthesis and they
determine the datapath contained in the bitstream. The network parameters --- weights,
bias, activation and quantization parameters, specific constants --- are instead loaded
at runtime through the host interface and stored in the RAM associated with the FPGA.

\begin{fnnote}[Central architectural principle]
A build fixes the \emph{ceiling} of the machine (maximum number of layers, maximum
width, \code{PARALLEL}); the host configures the \emph{actual} network --- number of
layers, per-layer input/output width, per-layer activation and trained parameters ---
entirely at runtime, over SPI, into the FPGA's local memory. A single bitstream serves
any topology up to that ceiling.
\end{fnnote}

\section{Boot and initialization}
The FPGA is configured at power-on through the usual configuration mechanism (bitstream
loading from SPI flash). The bitstream defines the hardware architecture of the engine;
the host does not dynamically build the datapath during normal operation, but rather
configures the network data on which the already existing datapath operates.

\begin{center}
\begin{tikzpicture}[font=\scriptsize,node distance=4.5mm,start chain=going below,
                    every node/.style={on chain}]
  \node[fnblockA,minimum width=60mm](p){Power-on};
  \node[fnblock,minimum width=60mm]{FPGA configuration (bitstream from flash)};
  \node[fnblockT,minimum width=60mm]{Neural Network Engine available};
  \node[fnblock,minimum width=60mm]{Host initialization (SPI)};
  \node[fnblock,minimum width=60mm]{Loading network parameters / weights / bias};
  \node[fnblockD,minimum width=60mm]{Engine ready};
  \begin{scope}[every path/.style={fnarrow}]
   \foreach \a/\b in {1/2,2/3,3/4,4/5,5/6}{}
  \end{scope}
  \foreach \i [count=\j from 2] in {1,...,5}{
     \draw[fnarrow] (chain-\i) -- (chain-\j);}
\end{tikzpicture}
\end{center}

\section{Training and inference}
Training and inference are conceptually separate. The first implementation does not
require the FPGA to perform training: weights can be computed externally
(PC/Linux/other host) and transferred over SPI into the FPGA's RAM, which then performs
inference. This drastically reduces the complexity of the initial hardware, without
precluding a future implementation of assisted or fully hardware training (roadmap
Phase~8, ch.~\ref{ch:roadmap}). During inference the host only provides the input data
and retrieves the result, obtaining deterministic computation, reduced host load,
hardware parallelism, predictable latency and independence from the host CPU
architecture.

\section{Design philosophy and reuse}
The project should be understood as a \emph{reusable FPGA neural acceleration platform}
rather than a single network. The application determines input size, topology, number
of layers and neurons, parallelism, numeric precision, activation functions, memory and
performance requirements; the hardware generation process produces the corresponding
FPGA implementation. The same HDL architecture remains conceptually unchanged while the
synthesis parameters generate implementations appropriate to the different application
targets.
